% ----------------------------------------------------------------------
%  3.5.4. Электросопротивление и строение сплавов; 3.5.5. Холодная
%  деформация; 3.5.6. Фазовые превращения
% ----------------------------------------------------------------------
\subsection{Связь электросопротивления со строением сплавов}\label{ssec:rho-alloys}

\paragraph{Твёрдые растворы}
Растворение второго компонента \emph{всегда} увеличивает сопротивление~—
даже если сам добавленный металл проводит лучше основы. Главная
причина~— искажения решётки, рассеивающие электроны. Эффект бывает
большим: 0,2~ат.\,\% As или Fe удваивают сопротивление золота.

Вторая причина~— химическое различие компонентов: чем больше разница
валентностей, тем больше добавочное сопротивление. Для малых
концентраций непереходных металлов выполняется \term{правило Линде}:
\begin{equation}\label{eq:linde}
\rho=\rho_0+p\,\bigl[a+b\,(z-z_\text{р})^2\bigr],
\end{equation}
где $p$~— концентрация добавки (ат.\,\%), $z$ и $z_\text{р}$~— валентности
добавки и растворителя, $a,b\approx0{,}3$--$0{,}5$~мкОм$\cdot$см на 1~ат.\,\%\ispr{В
оригинале (с.~507) размерность констант указана как «Ом$\cdot$см»;
добавочное сопротивление в правиле Линде измеряется в мкОм$\cdot$см на
1~ат.\,\% добавки.}.

\paragraph{Как сопротивление зависит от состава}
Вид кривой «сопротивление~— состав» определяется типом диаграммы
состояния (рис.~\ref{fig:kurnakov}; эти закономерности установил
Н.\,С.~Курнаков).
\begin{itemize}
\item \emph{Непрерывный ряд твёрдых растворов} близких по строению
металлов (Cu--Ag, Cu--Au, Ag--Au): кривая~— парабола с максимумом при
равном числе атомов, где решётка искажена сильнее всего
(рис.~\ref{fig:ss-props}). Если же электронное строение компонентов
различно (благородный металл~— переходный), максимум сдвигается к
переходному металлу: валентные электроны переходят на его
недостроенные $d$-уровни, и электронов проводимости становится меньше.
\item \emph{Упорядочение} делает поле решётки симметричнее, длина
пробега растёт, сопротивление падает (Cu--Au, рис.~\ref{fig:CuAu-rho}).
По тому, насколько близко к прямой, соединяющей чистые металлы, опускаются
минимумы при Cu$_3$Au и CuAu, судят о степени дальнего порядка. Упорядочение
возможно и в промежуточных фазах и тоже снижает сопротивление.
\item \emph{Гетерогенные (двухфазные) сплавы}: проводимость
складывается из проводимостей фаз и меняется с объёмной долей фаз
линейно~— если кристаллы фаз крупные, равноосные, а проводимости фаз
различаются не сильно. Отклонения возникают, когда частицы одной из фаз
соизмеримы с длиной свободного пробега: на них электроны рассеиваются
дополнительно (максимум эффекта~— при размерах около 1~нм, рост
сопротивления до 10--15\,\%). Отклоняются от аддитивности и многие
эвтектики (искажения на развитых границах).
\item \emph{Химические соединения} проводят, как правило, гораздо хуже
компонентов: металлическая связь частично заменяется ковалентной или
ионной, и носителей становится меньше. Некоторые соединения~— даже
полупроводники (Mg$_2$Sn, ZnSb, Mg$_3$Sb$_2$); у них при отклонении от
стехиометрии сопротивление резко падает. Для дальтонидов экстремум
свойства лежит внутри области гомогенности фазы, для бертоллидов~— на
её границе. Металлическую проводимость имеют алюминиды, нитриды,
бориды, карбиды; силициды (особенно богатые кремнием) часто~—
полупроводники.
\end{itemize}

\begin{figure}[ht]
\centering
\begin{tikzpicture}
\begin{groupplot}[group style={group size=4 by 2, horizontal sep=3mm, vertical sep=3mm,
    xticklabels at=edge bottom, yticklabels at=edge left},
  width=0.31\textwidth, height=3.1cm, xmin=0, xmax=1, xtick=\empty, ytick=\empty, axis lines=box,
  every axis/.append style={font=\scriptsize}]
  % --- диаграммы состояния
  \nextgroupplot[ymin=0, ymax=1, title={\textit{а}}]
    \addplot[atomLine, smooth, domain=0:1] {0.85-0.3*x+0.12*sin(deg(pi*x))};
    \addplot[atomLine, smooth, domain=0:1] {0.85-0.3*x-0.08*sin(deg(pi*x))};
    \node at (axis cs:0.5,0.3) {$\alpha$};
  \nextgroupplot[ymin=0, ymax=1, title={\textit{б}}]
    \addplot[atomLine] coordinates {(0,0.85) (0.45,0.45) (1,0.8)};
    \addplot[atomLine] coordinates {(0,0.45) (1,0.45)};
    \node at (axis cs:0.5,0.22) {A + B};
  \nextgroupplot[ymin=0, ymax=1, title={\textit{в}}]
    \addplot[atomLine] coordinates {(0,0.85) (0.45,0.45) (1,0.8)};
    \addplot[atomLine] coordinates {(0.15,0.45) (0.8,0.45)};
    \addplot[atomLine, smooth] coordinates {(0.15,0.45) (0.12,0.25) (0.1,0)};
    \addplot[atomLine, smooth] coordinates {(0.8,0.45) (0.85,0.25) (0.88,0)};
    \node at (axis cs:0.47,0.22) {$\alpha+\beta$};
  \nextgroupplot[ymin=0, ymax=1, title={\textit{г}}]
    \addplot[atomLine] coordinates {(0,0.75) (0.22,0.45) (0.5,0.9) (0.75,0.5) (1,0.7)};
    \addplot[atomLine] coordinates {(0,0.45) (0.5,0.45)};
    \addplot[atomLine] coordinates {(0.5,0.5) (1,0.5)};
    \addplot[atomLine] coordinates {(0.5,0) (0.5,0.9)};
    \node at (axis cs:0.5,0.95) {\tiny A$_m$B$_n$};
  % --- сопротивление
  \nextgroupplot[ymin=0, ymax=1.1, ylabel={$\rho$}]
    \addplot[fmRed, thick, domain=0:1, samples=40] {0.15+0.1*x+3.2*x*(1-x)};
  \nextgroupplot[ymin=0, ymax=1.1]
    \addplot[fmRed, thick] coordinates {(0,0.15) (1,0.3)};
  \nextgroupplot[ymin=0, ymax=1.1]
    \addplot[fmRed, thick, smooth, domain=0:0.1, samples=10] {0.15+5*x};
    \addplot[fmRed, thick] coordinates {(0.1,0.65) (0.88,0.75)};
    \addplot[fmRed, thick, smooth, domain=0.88:1, samples=10] {0.75-3.75*(x-0.88)};
  \nextgroupplot[ymin=0, ymax=1.1]
    \addplot[fmRed, thick] coordinates {(0,0.15) (0.5,0.95) (1,0.3)};
\end{groupplot}
\end{tikzpicture}
\caption{Диаграммы состояния (вверху) и соответствующие зависимости
сопротивления от состава (внизу), схема по Курнакову: \textit{а}~—
непрерывный ряд твёрдых растворов (парабола); \textit{б}~— механическая
смесь компонентов (прямая); \textit{в}~— ограниченные растворы и
двухфазная область (подъём в однофазных областях, прямая в двухфазной);
\textit{г}~— химическое соединение (особая точка на кривой при составе
соединения)}
\label{fig:kurnakov}
\end{figure}

\subsection{Влияние холодной деформации на электросопротивление}\label{ssec:rho-cold}

Сопротивление чистых металлов даже при сильном наклёпе растёт всего на
2--6\,\%: сказываются искажения решётки. Наклёпанный металл имеет
повышенное остаточное сопротивление; отжиг частично устраняет дефекты,
и после рекристаллизации сопротивление возвращается к исходному. Так
же ведут себя однородные твёрдые растворы.

Есть и обратные случаи.
\begin{itemize}
\item Сплавы с \emph{неоднородными} твёрдыми растворами (Ni--Cr, Ni--Cu,
Fe--Al, Cu--Mn, Au--Cr): деформация разрушает неоднородность и
расставляет атомы статистически, и сопротивление \emph{падает}, а при
рекристаллизационном отжиге снова растёт.
\item \emph{Гетерогенные} сплавы: важно не только число дефектов, но и
взаимное расположение фаз. Сопротивление стальной проволоки с 0,58\,\%
углерода при волочении \emph{уменьшается} (на 5,5\,\%): пластинки
цементита (с высоким сопротивлением) ориентируются вдоль проволоки, и
эффект ориентировки превышает рост сопротивления наклёпанного феррита.
Отпуск при 400--650\,°C снижает сопротивление ещё больше, а отжиг при
более высоких температурах уничтожает ориентировку и повышает его.
\end{itemize}

\subsection{Изменение электрических свойств при фазовых превращениях}\label{ssec:el-phase}

Зависимость сопротивления от фазового и структурного состояния
позволяет изучать превращения, просто измеряя сопротивление.
\begin{itemize}
\item \emph{Магнитные превращения.} У ферромагнетиков ниже точки Кюри
сопротивление падает быстрее, чем у сходного парамагнитного металла
(рис.~\ref{fig:Ni-Pd}): с ростом спонтанного порядка спинов исчезает
рассеяние электронов на беспорядке спинов.
\item \emph{Плавление.} Сопротивление жидкого металла в 1,5--2 раза больше,
чем твёрдого (нарушение порядка, изменение связи). Для большинства
металлов (кроме Ga, Hg, Sb, Bi)
\begin{equation}\label{eq:mott-melt}
\left(\frac{\rho_\text{ж}}{\rho_\text{тв}}\right)_{\Tpl}=\exp\!\left(\frac{80\,L_\text{пл}}{\Tpl}\right),
\end{equation}
где $L_\text{пл}$~— теплота плавления в кДж/моль. Например, для меди
($L_\text{пл}=13{,}1$~кДж/моль, $\Tpl=1358$~К) получаем 2,2 при опытном
значении около 2,1. У сурьмы сопротивление при плавлении, наоборот,
\emph{падает}: в твёрдом состоянии у неё заметна ковалентная связь, а в
жидком связь становится металлической.
\item \emph{Диаграммы состояния и распад.} Измерением сопротивления
определяют линии ограниченной растворимости, изучают изотермический
распад аустенита (начало и конец превращения электрическим методом
фиксируются точнее, чем магнитным), закалку, отпуск, дисперсионное
твердение.
\item \emph{Старение} (рис.~\ref{fig:aging}). У сплава Al~+~4,5\,\% Cu при
20\,°C сопротивление \emph{растёт}: образуются зоны Гинье--Престона,
сильно рассеивающие электроны. При 225\,°C оно \emph{падает}: из раствора
выделяется соединение CuAl$_2$ со своей решёткой, и раствор обедняется.
\end{itemize}

\begin{figure}[ht]
\centering
\begin{subfigure}[b]{0.48\textwidth}\centering
\begin{tikzpicture}
\begin{axis}[width=\textwidth, height=4.8cm, xmin=0, xmax=1.3, ymin=0, ymax=1.35,
  xlabel={$T/\theta_K$ (Ni)}, ylabel={$\rho/\rho_\theta$}, axis lines=left, xtick={0,0.5,1}, ytick={0,0.5,1}]
  \addplot[fmBlue, thick, domain=0:1.3, samples=60] {x<1 ? 0.03+0.97*x^1.8 : 1+0.35*(x-1)};
  \addplot[fmRed, thick, domain=0:1.3, samples=40] {0.05+0.85*x-0.05*x^2};
  \draw[fmGray, dotted] (axis cs:1,0) -- (axis cs:1,1.35);
  \node[font=\scriptsize, fmBlue, anchor=north west] at (axis cs:0.55,0.28) {Ni (ферро)};
  \node[font=\scriptsize, fmRed, anchor=south east] at (axis cs:0.5,0.55) {Pd (пара)};
\end{axis}
\end{tikzpicture}
\caption{}
\label{fig:Ni-Pd}
\end{subfigure}\hfill
\begin{subfigure}[b]{0.48\textwidth}\centering
\begin{tikzpicture}
\begin{axis}[width=\textwidth, height=4.8cm, xmin=0, xmax=10, ymin=-1, ymax=1,
  xlabel={время старения}, ylabel={$\Delta\rho$}, xtick=\empty, ytick={0}, axis lines=left]
  \draw[fmGray] (axis cs:0,0) -- (axis cs:10,0);
  \addplot[fmBlue, thick, domain=0:10, samples=60] {0.7*(1-exp(-x/2))};
  \addplot[fmRed, thick, domain=0:10, samples=60] {0.35*exp(-x/0.6)*x/0.6*2.7/2.7-0.85*(1-exp(-x/3))};
  \node[font=\scriptsize, fmBlue, anchor=south east] at (axis cs:9.8,0.72) {20\,°C: зоны ГП};
  \node[font=\scriptsize, fmRed, anchor=north east] at (axis cs:9.8,-0.75) {225\,°C: CuAl$_2$};
\end{axis}
\end{tikzpicture}
\caption{}
\label{fig:aging}
\end{subfigure}
\caption{\textit{а}~— температурная зависимость сопротивления
ферромагнитного никеля и парамагнитного палладия, отнесённого к
сопротивлению при точке Кюри никеля (схема): ниже $\theta_K$ у никеля
сопротивление падает резче; \textit{б}~— изменение сопротивления при
старении сплава Al~+~4,5\,\% Cu (схема)}
\end{figure}
